\chapter{Core Values and Observable State}
\label{ch:core-values}
\chapterquote{One changing fact deserves one owner and one observable contract.}{The value principle}

\begin{chaptermap}
Core question & How does Corusco represent one changing screen fact without depending on Swing?\\
Primary APIs & \api{ReadableValue}, \api{WritableValue}, \api{SimpleValue}, \api{MappedValue}, \api{DerivedValue}, \api{MasterDetailValue}, \api{LoadableValue}, \api{ValueChangeEvent}, and \api{ChangeOrigin}.\\
Plain Swing baseline & Labels, buttons, selections, and fields often begin as component state plus listeners; the chapter separates the underlying facts from the widgets.\\
Swing connection & Swing bindings mirror values into labels, fields, status text, selection state, commands, busy indicators, and overlays.\\
Companion example & \coreExample.\\
\end{chaptermap}

\section{A value is a small model}

\termdef{Observable value}{in Corusco} is a scalar model with a current value and
a subscription contract. It is smaller than a form model, smaller than a table
model, and larger than a variable. A variable can be read and written; an
observable value can also tell interested code when the effective value changes.

This is the first important Corusco idea because many Swing screens begin by
letting widgets own facts. A \api{JLabel} owns the current status text. A
\api{JButton} owns whether Save is enabled. A \api{JTable} owns the selected row.
A \api{JTextField} owns the current input. This works for a demonstration. It
does not scale when the same fact must appear in a toolbar, menu, dialog
summary, keyboard shortcut, status line, test, and generated binding.

The value contract moves the fact out of the widget:

\begin{lstlisting}[caption={One changing fact with an explicit owner.}]
WritableValue<String> selectedCustomerName = SimpleValue.of("Ada Lovelace");

Subscription subscription = selectedCustomerName.subscribe(event -> {
    System.out.println(event.oldValue() + " -> " + event.newValue());
});

selectedCustomerName.setValue("Grace Hopper", ChangeOrigin.USER);
subscription.close();
\end{lstlisting}

Nothing in this listing depends on Swing. The same value could later feed a
label, a window title, a command, or a test assertion. The widget is no longer
the only place where the fact exists.

\begin{principlebox}{small models first}
Before inventing a presenter hierarchy, ask whether the screen fact is simply
one observable value. A selected row, busy flag, disabled reason, active tab,
status message, or committed filter text often is.
\end{principlebox}

\textbf{Readable and writable contracts.}

\api{ReadableValue<T>} has two obligations. It returns the current value through
\api{value()}, and it accepts subscribers through \api{subscribe}. The
subscription returns a \api{Subscription}, which is the caller's cleanup handle.
\api{WritableValue<T>} adds \api{setValue}. The default \api{setValue(value)}
uses \api{ChangeOrigin.MODEL}; callers may supply another origin when the
difference matters.

\begin{lstlisting}[caption={The shape of the core contract.}]
public interface ReadableValue<T> {
    T value();
    Subscription subscribe(ValueChangeListener<T> listener);
}

public interface WritableValue<T> extends ReadableValue<T> {
    default void setValue(T value) {
        setValue(value, ChangeOrigin.MODEL);
    }

    void setValue(T value, ChangeOrigin origin);
}
\end{lstlisting}

The contract is deliberately small. It does not say "EDT". It does not say
"background thread". It does not say "component". Built-in values deliver events
synchronously on the thread that performs the mutation, and equal effective
values do not produce events. That policy is simple enough to test and simple
enough to reason about.

\begin{pitfallbox}{values are not hidden dispatchers}
Do not assume that a Corusco core value will marshal events to Swing for you.
Thread ownership belongs at the Swing boundary. If a value is consumed by Swing,
mutate it on the EDT or adapt it before Swing observes it.
\end{pitfallbox}

\textbf{Effective changes.}

A value change event represents an effective change, not merely a call to a
setter. Built-in implementations compare old and new values with
\api{Objects.equals}. Setting the same value again is silent. This is an
important difference from many hand-written listener systems, where every setter
call fires and listeners must defend themselves from redundant work.

\begin{lstlisting}[caption={Repeated equal values do not create useful events.}]
SimpleValue<String> mode = SimpleValue.of("view");
List<String> events = new ArrayList<>();

mode.subscribe(event -> events.add(event.oldValue() + ":" + event.newValue()));
mode.setValue("edit");
mode.setValue("edit");
mode.setValue("view");

// events contains "view:edit" and "edit:view"
\end{lstlisting}

This policy keeps derived values, bindings, and status updates from repainting
for no semantic reason. It also means that "refresh even if equal" is a
different operation. If a screen needs an explicit refresh signal, model that as
a command, generation value, or event source rather than abusing a scalar value.

\textbf{The event object.}

\api{ValueChangeEvent<T>} carries four pieces of information: the source value,
the old value, the new value, and the \api{ChangeOrigin}. Old and new values may
be null. The source and origin are never null.

The old value is useful for differential updates, undo stacks, and diagnostics.
The new value is the current fact. The source is useful when a listener observes
more than one value. The origin is the part that ordinary Swing listeners do not
have.

\begin{lstlisting}[caption={Using old value, new value, source, and origin.}]
ValueChangeListener<CustomerRow> listener = event -> {
    if (event.origin() == ChangeOrigin.USER) {
        auditSelection(event.oldValue(), event.newValue());
    }
    refreshStatusFor(event.source().value());
};
\end{lstlisting}

Do not overuse the source. Most code should subscribe to one value and already
know what it observes. But the source is part of the contract because it makes
composition possible without losing identity.

\textbf{Change origins.}

\termdef{Change origin}{in Corusco} is a small value describing why a mutation
occurred. Built-in origins include \api{USER}, \api{MODEL}, \api{BINDING},
\api{GENERATED}, and \api{SYSTEM}. Applications may create additional origins
with \api{ChangeOrigin.of}.

Origins are not a security system and not a replacement for commands. They are a
diagnostic and feedback-control tool. They answer questions that listeners
otherwise guess:

\begin{itemize}
\item Did this update come from the user or from loading a record?
\item Did a generated binding write the value while synchronizing a component?
\item Did a system reset clear the field?
\item Should this listener write back, or would that create a feedback loop?
\end{itemize}

\begin{lstlisting}[caption={Distinguishing user and model changes.}]
selectedCustomer.subscribe(event -> {
    if (event.origin() == ChangeOrigin.USER) {
        rememberRecentSelection(event.newValue());
    }
});

selectedCustomer.setValue(rowFromTable, ChangeOrigin.USER);
selectedCustomer.setValue(rowFromInitialLoad, ChangeOrigin.MODEL);
\end{lstlisting}

In a table selection binding, this distinction matters. A user click should be
reported as user-originated. A presenter selecting a row after loading data is
model-originated. Both changes may lead to the same selected row, but they have
different meaning.

\begin{coruscobox}{origins at the Swing boundary}
Corusco Swing bindings write user-originated component edits with
\api{ChangeOrigin.USER}. Presenter-originated changes retain their model origin
when they update components. This makes feedback loops easier to detect and
keeps tests from asserting only the final value.
\end{coruscobox}

\textbf{Subscriptions are resources.}

Subscribing installs a listener. A listener is a resource. A resource needs an
owner. This is the most mundane rule in the chapter, and one of the most
important.

Swing programs often leak because the visual component disappears but a listener
or task callback still references it. Corusco makes the cleanup handle visible:
\api{subscribe} returns a \api{Subscription}. Closing it removes the listener.
The close operation is idempotent.

\begin{lstlisting}[caption={A screen owns the subscriptions it creates.}]
final class CustomerStatusPresenter implements AutoCloseable {
    private final List<Subscription> subscriptions = new ArrayList<>();

    CustomerStatusPresenter(ReadableValue<CustomerRow> selected,
            WritableValue<String> statusText) {
        subscriptions.add(selected.subscribe(event -> {
            CustomerRow row = event.newValue();
            statusText.setValue(row == null ? "No customer" : row.displayName());
        }));
    }

    @Override
    public void close() {
        subscriptions.forEach(Subscription::close);
        subscriptions.clear();
    }
}
\end{lstlisting}

Later chapters use scopes to group subscriptions and Swing bindings. The core
lesson is already visible: if you install a relationship, you must own its
removal.

\section{SimpleValue}

\api{SimpleValue<T>} is the ordinary mutable implementation. It stores one
current value, accepts subscribers, compares effective changes, and emits
synchronous events. Use it for presentation facts owned by a presenter or model.
Do not use it as a global variable with better manners.

\begin{lstlisting}[caption={Simple presentation facts.}]
SimpleValue<Boolean> busy = SimpleValue.of(false);
SimpleValue<String> statusText = SimpleValue.of("Ready");
SimpleValue<CustomerRow> selectedCustomer = SimpleValue.empty();
\end{lstlisting}

The important design question is not "can this be a \api{SimpleValue}?" It
usually can. The question is "who owns this value?" A screen presenter may own a
busy flag. A long-lived application model may own the active workspace. A dialog
form may own dirty state. Ownership should follow lifetime.

\begin{pitfallbox}{observable globals}
Replacing global mutable variables with global observable values does not fix
global state. It merely makes the global state louder. Prefer values whose
lifetime matches the feature they represent.
\end{pitfallbox}

\textbf{MappedValue.}

\api{MappedValue<A,B>} is the one-source derived case. It observes a source
value and exposes a mapped value. This is useful for labels, command enablement,
tooltip fragments, and small view-state conversions.

\begin{lstlisting}[caption={A display value derived from a selected row.}]
ReadableValue<CustomerRow> selected = presenter.selectedCustomer();

MappedValue<CustomerRow, String> title = MappedValue.of(selected, row ->
        row == null ? "Customers" : "Customer: " + row.displayName());
\end{lstlisting}

Without a mapped value, the code is usually written as a listener that updates a
nearby variable or component. The dependency then hides inside the listener. With
a mapped value, the dependency is itself a model object and can be subscribed to,
tested, and closed.

\textbf{DerivedValue.}

\api{DerivedValue<T>} generalizes mapping to multiple dependencies. Use it when
the output depends on more than one value: Save enablement depends on dirty
state, validation problems, and busy state; status text may depend on selection
and load state; a dialog title may depend on mode and current record.

\begin{lstlisting}[caption={Derived command enablement.}]
DerivedValue<Boolean> canSave = DerivedValue.of(
        () -> form.dirty().value()
                && !form.problems().value().hasErrors()
                && !busy.value(),
        form.dirty(),
        form.problems(),
        busy);
\end{lstlisting}

This listing also shows a useful test shape. A test can change dirty state,
problems, and busy state, then assert \api{canSave.value()} and its emitted
events. No button is required. The button binding is a later concern.

\begin{principlebox}{derived before listener}
When a listener merely recomputes a value from other values, consider a derived
value. The dependency graph becomes visible, and the listener body disappears.
\end{principlebox}

\textbf{Avoiding derived-value soup.}

Derived values can also be overused. A screen with fifty tiny derived values and
no names for higher-level concepts may be as hard to read as a screen with fifty
listeners. Group related facts into a presenter or form model when the
relationships become a vocabulary of their own.

The test is language. If you can name the derived value as a domain of the
screen, keep it. \api{canSave}, \api{primaryStatusText}, and
\api{selectedCustomerTitle} are good names. If the name becomes
\api{labelTextWhenDirtyAndLoadedButNotSaving}, the concept may need a state
object instead.

\textbf{Master-detail values.}

Master-detail screens are common in Swing. Selecting a customer loads orders.
Selecting a project loads tasks. Selecting a file loads metadata. Hand-written
Swing code often puts this relationship inside a table selection listener. The
listener reads the selected row, calls a loader, updates detail fields, handles
empty selection, and remembers cleanup. That is too much work for a listener.

\api{MasterDetailValue<M,T>} models the relationship between a master value and
a detail value. The master is observable. The detail is loaded from the master.
When the master changes, the detail updates.

\begin{lstlisting}[caption={Detail value follows the selected master.}]
SimpleValue<CustomerId> selectedCustomerId = SimpleValue.empty();

MasterDetailValue<CustomerId, CustomerDetails> details =
        MasterDetailValue.of(selectedCustomerId, repository::loadDetails);

selectedCustomerId.setValue(new CustomerId(42), ChangeOrigin.USER);
\end{lstlisting}

The example is synchronous because core values are synchronous. If detail
loading may block, move the loading into a task chapter pattern or use a
loadable/task boundary. The conceptual relationship is still master-detail:
selection determines detail.

\textbf{Loadable values.}

\api{LoadableValue<T>} represents a lazily or explicitly loaded scalar value.
It is useful when constructing the screen should not immediately construct the
expensive value, or when a cached value can be detached and later reloaded. A
loadable value keeps the loading policy close to the value rather than scattered
through listeners.

\begin{lstlisting}[caption={A cached scalar value with explicit loading.}]
AtomicInteger loads = new AtomicInteger();
LoadableValue<String> customerName =
        LoadableValue.of(() -> "customer-" + loads.incrementAndGet());

String first = customerName.value();
String second = customerName.value(); // Cached according to the value contract.
customerName.detach();
String third = customerName.value();  // Loaded again.
\end{lstlisting}

The exact caching and detaching behavior belongs to the implementation contract.
The design idea is simple: loading is not an accidental side effect hidden in a
label update. It is part of the value model.

\section{Null as a state}

Corusco values may contain null. That is a practical choice because Swing and
desktop workflows often have absent selection, empty optional fields, and
not-yet-loaded data. But null should still be designed. A null selected row is
reasonable. A null status text is usually poor design. A null validation problem
set is unnecessary when an empty \api{ProblemSet} exists.

When the set of states grows beyond "present or absent", prefer a state object:
\api{Idle}, \api{Loading}, \api{Ready}, and \api{Failed} are clearer than a
value plus several flags plus null.

\textbf{Plain Swing baseline.}

The plain Swing version of a status update usually looks like this:

\begin{lstlisting}[caption={Component-owned status.}]
table.getSelectionModel().addListSelectionListener(event -> {
    int viewRow = table.getSelectedRow();
    if (viewRow < 0) {
        statusLabel.setText("No customer selected");
        editButton.setEnabled(false);
        return;
    }
    int modelRow = table.convertRowIndexToModel(viewRow);
    CustomerRow row = model.rowAt(modelRow);
    statusLabel.setText(row.displayName());
    editButton.setEnabled(true);
});
\end{lstlisting}

This is not terrible code. It may be exactly right for a small utility panel.
The problem appears when menu items, detail panes, tests, and persistence also
need selected-customer state. Then the status label and button listener are the
wrong owners.

\textbf{Corusco treatment.}

The Corusco version names the selected row, derives the status, and binds the
component later.

\begin{lstlisting}[caption={Value-owned status.}]
WritableValue<CustomerRow> selectedCustomer = SimpleValue.empty();

ReadableValue<String> statusText = MappedValue.of(selectedCustomer, row ->
        row == null ? "No customer selected" : row.displayName());

ReadableValue<Boolean> canEdit = MappedValue.of(selectedCustomer,
        Objects::nonNull);
\end{lstlisting}

The table selection binding writes \api{selectedCustomer}. The label binding
reads \api{statusText}. The button or command binding reads \api{canEdit}. Each
relationship can be closed when the screen closes. Tests can assert the model
without constructing the table.

\textbf{Testing value state.}

Core values are intentionally easy to test. A good value test does not click a
button merely to prove that a derived value works. It mutates the source values,
records events, and asserts final state and origins.

\begin{lstlisting}[caption={Testing derived state directly.}]
SimpleValue<Boolean> dirty = SimpleValue.of(false);
SimpleValue<Boolean> busy = SimpleValue.of(false);
SimpleValue<ProblemSet> problems = SimpleValue.of(ProblemSet.empty());

DerivedValue<Boolean> canSave = DerivedValue.of(
        () -> dirty.value() && !problems.value().hasErrors() && !busy.value(),
        dirty, problems, busy);

dirty.setValue(true, ChangeOrigin.USER);
assertThat(canSave.value()).isTrue();

busy.setValue(true, ChangeOrigin.SYSTEM);
assertThat(canSave.value()).isFalse();
\end{lstlisting}

The Swing binding should have its own small test, but it should not be the only
test of the state rule.

\textbf{When a value is not enough.}

Use a value for one changing fact. Use a form model for a group of editable
fields with parsing, validation, dirty state, and commit/reset semantics. Use an
observable list for a changing collection. Use a command for an operation with
enabled state, selected state, resources, and invocation. Use a task handle for
background work. A value is the atom of presentation state, not the whole
periodic table.

\begin{migrationbox}{from listener variables to values}
During migration, look for private fields that are changed by listeners and read
by other listeners. If the field represents one screen fact, replace it with a
\api{SimpleValue}. If it represents several facts, introduce a named model
instead.
\end{migrationbox}

\textbf{Worked example.}

\begin{examplebox}{core values}
\coreExample{} creates selected-customer state, derived status text, and scoped
cleanup. The example is deliberately Swing-free so the state contract can be
tested directly before a component binding is introduced.
\end{examplebox}

Read the example by following ownership. Which object owns the source value?
Which object owns the derived value? Which code owns the subscription? Which
change is user-originated and which is model-originated? If those answers are
clear, the eventual Swing code will have fewer surprises.

\textbf{Review notes.}

Values make state visible, but they do not make design automatic. A value with a
bad name is still a bad name. A value with the wrong lifetime still leaks. A
value mutated from the wrong thread still violates a Swing boundary if Swing is
observing it. The contract is intentionally small so ownership remains visible.

The first useful question is whether the fact has a single owner. A selected
customer, current filter text, busy flag, status message, or active problem set
often does. A complex form, editable table, or workflow state machine often
does not. If a value begins to carry several meanings in its name, such as
\api{state} or \api{modeAndSelection}, it is no longer helping the reader. Split
the fact or introduce a model with named members.

The second question is lifetime. A value owned by a presenter may outlive the
Swing component that displays it. A value owned by a dialog should normally die
with the dialog. A global preference value may live for the application. The
value API deliberately does not guess these lifetimes. Subscriptions and
derived values must be registered with the object that knows when the
relationship ends.

Origins are useful when a value participates in feedback. A text field update
from the user and a model reset from the application may produce the same new
string, but they do not have the same meaning. A binding may suppress a loop
when it sees that the model change came from the component it just updated. A
test may assert that accepting a baseline produces a model-originated change
rather than pretending it was user input.

Derived values should be pure in spirit even when their implementation is
ordinary Java. They should read source values and compute a result. They should
not mutate components, start tasks, or write back to their sources. Once a
derived value has side effects, a reader cannot reason about it as state
anymore. Use a subscription or command handler for effects; use a derived value
for a fact.

Null deserves an explicit policy. Sometimes absence is naturally represented by
\api{null}, as with no selected row in a simple screen. Sometimes a domain type
should use an optional-like wrapper or a sealed hierarchy to distinguish absent,
loading, failed, and available states. The important part is consistency. A
screen where some values use null, some use empty strings, and some use sentinel
objects forces every binding to rediscover absence rules.

Mapped values are most helpful when they turn a domain-shaped fact into a
presentation-shaped fact. A selected row becomes a status string. A problem set
becomes a boolean committable flag. A task handle becomes an enabled state. The
mapping is not a place to hide policy that belongs elsewhere. If computing the
mapping requires several domain services or user permissions, the presenter
probably needs a named method or model rather than an inline lambda.

Master-detail values are another ownership test. The master value says which
object is current. The detail value loads or derives data for that object. A
late detail result for an old master must not replace the current detail. This
is the same stale-result problem seen in background tasks, only expressed as
state. The value abstraction helps because the master relationship is explicit
and can be closed or invalidated.

Loadable values should be honest about their states. "No value yet" is not the
same as "loading", "failed", or "loaded empty". A screen that collapses all of
those into null cannot explain itself to the user. The view may display a simple
message, but the model should know which state it is in. That distinction also
makes tests clearer: one can assert that retry changes failed to loading, then
loading to loaded.

Threading policy should remain at the boundary. Core values are useful outside
Swing and should not smuggle Swing assumptions into their center. If a value is
observed by Swing, adapt or subscribe at the EDT boundary. If a value is updated
from background work, decide where the handoff occurs. This keeps the core
model testable and lets non-Swing tests run without an event queue.

Values are often the first Corusco object introduced during migration because
they have low ceremony. That does not mean every field in a class should become
a value immediately. Start with facts that are read in more than one place,
facts whose changes should be tested, and facts whose current listener owner is
clearly wrong. Leave local temporary variables alone. Good migration makes
state clearer; it does not turn every line into framework vocabulary.

The payoff appears in code review. Instead of searching through listeners to
discover why Save is enabled, the reviewer can inspect a derived
\api{canSave}. Instead of asking whether a closed panel still receives updates,
the reviewer can find the scope that owns the subscription. Instead of clicking
a UI to prove a status message, a test can update the selected value and assert
the mapped text. These small gains compound across a Swing application.

\section{Designing Value Graphs in Real Screens}

A value by itself is simple. A screen made of values is not automatically simple.
The useful design object is the \termdef{value graph}{the small directed graph of
source values, derived values, subscriptions, bindings, and command facts owned
by one presenter or workflow}. In a good value graph, facts flow from owners to
observers. Source values receive user or model mutations. Derived values compute
secondary facts. Bindings and command adapters present those facts to Swing.
Effects happen at explicit subscription or command boundaries. The graph can be
read without opening a debugger.

This vocabulary is deliberately modest. Corusco values are not a spreadsheet, not
a transaction engine, and not a hidden reactive runtime. They are small
synchronous model objects that make dependency edges visible. That is enough to
remove much of the uncertainty from Swing code. A reviewer can ask: what is the
source of this fact, what values depend on it, who owns the subscriptions, and
where do effects occur? If the answer is spread across anonymous listeners on
five components, the screen is still under-modeled.

Start a value graph by naming source facts. A source fact is a fact that some
owner is allowed to write directly: selected customer, current filter text,
dirty flag, active tab, busy flag, current problem set, current load state, or
chosen export target. A source fact should have one clear owner. A text binding
may write \api{filterText} with a user origin. A repository load may write
\api{loadState} with a model or system origin. A dialog controller may write
\api{dirty}. Several components may observe these values, but the writing policy
should not be scattered.

Derived facts come next. A derived fact answers a question from source facts:
can the Save command run, what should the status line say, is the detail pane
visible, should the table show the empty-state overlay, what tooltip should a
disabled command show? When a listener only recomputes such a fact, make the
fact explicit. The result is not only less code; it is a better sentence. The
screen no longer says "when this checkbox changes, update these three widgets."
It says "committable is dirty and valid and not busy."

\begin{lstlisting}[caption={One presenter value graph in ordinary Java.}]
final class CustomerPresenter implements AutoCloseable {
    private final SimpleValue<CustomerRow> selected = SimpleValue.empty();
    private final SimpleValue<String> filterText = SimpleValue.of("");
    private final SimpleValue<Boolean> dirty = SimpleValue.of(false);
    private final SimpleValue<Boolean> busy = SimpleValue.of(false);
    private final SimpleValue<ProblemSet> problems =
            SimpleValue.of(ProblemSet.empty());

    private final DerivedValue<Boolean> canSave = DerivedValue.of(
            () -> dirty.value()
                    && !busy.value()
                    && !problems.value().hasErrors(),
            dirty, busy, problems);

    private final DerivedValue<String> statusText = DerivedValue.of(
            () -> selected.value() == null
                    ? "No customer selected"
                    : "Selected " + selected.value().displayName(),
            selected);

    @Override
    public void close() {
        canSave.close();
        statusText.close();
    }
}
\end{lstlisting}

The listing is intentionally not a framework showcase. It is a readable state
graph. A component layer can later bind a table selection to \api{selected}, a
text field to \api{filterText}, a button or command to \api{canSave}, and a
status label to \api{statusText}. The presenter can be tested without creating
any component. The values have names that a code reviewer can discuss.

A good source value has a strong name and a narrow type. \api{selectedCustomer}
is better than \api{selected}. \api{selectedCustomerId} may be better still when
the table row object is a view adapter and the durable identity is an id.
\api{loadState} is better than three booleans named \api{loading},
\api{loaded}, and \api{failed}. \api{currentMode} may be too vague unless the
mode type is a small explicit hierarchy. The chapter's first rule, one changing
fact per value, depends on the type being honest about the fact.

Records and sealed types help when a scalar value has several meaningful states.
A loading screen is rarely just null or not null. It may be idle, loading,
loaded with data, loaded empty, failed, refreshing existing data, or cancelled.
A single \api{ReadableValue<LoadState>} can express that vocabulary more clearly
than several booleans and a nullable exception. The value remains scalar, but the
state it carries is structured.

\begin{lstlisting}[caption={A scalar value can carry a structured state.}]
sealed interface CustomerLoadState {
    record Idle() implements CustomerLoadState {}
    record Loading(String query) implements CustomerLoadState {}
    record Loaded(List<CustomerRow> rows) implements CustomerLoadState {}
    record Failed(String query, Throwable cause) implements CustomerLoadState {}
}

SimpleValue<CustomerLoadState> customers =
        SimpleValue.of(new CustomerLoadState.Idle());
\end{lstlisting}

This is often better than making every detail independently observable. A
listener that observes \api{loading}, \api{rows}, \api{error}, and \api{query}
may see intermediate combinations that never make sense to the user. A single
state value says which combinations are legal. Derived values can still expose
presentation facts such as \api{isBusy}, \api{emptyMessage}, or
\api{retryEnabled}. The source remains one coherent fact.

Equality is part of the graph. Built-in values compare effective changes with
\api{Objects.equals}. Records are useful because their equality usually matches
their visible state. Mutable objects are dangerous when stored in values because
the value cannot see internal mutation. If a \api{List} is stored inside a
value and later modified in place, no setter call occurs and no value event is
published. For changing collections, use an observable list. For scalar values,
prefer immutable records, ids, enums, or small state objects.

\begin{pitfallbox}{mutable payloads in scalar values}
A \api{SimpleValue<List<Row>>} does not become an observable list. If callers
mutate the list in place, observers see no change event and derived values may
remain stale. Use immutable list snapshots for scalar state, or use the
observable-list chapter's collection model when the collection itself changes.
\end{pitfallbox}

Derived values should be pure in spirit. The supplier should read dependencies
and compute a result. It should not write a source value, mutate a component,
start a repository call, or show a dialog. This is not moralism; it prevents
cycles. If \api{canSave} recomputes and writes \api{dirty}, and \api{dirty}
causes \api{canSave} to recompute, the graph has become a hidden control flow
machine. Effects belong in command handlers or explicitly owned subscriptions.

Some cycles are legitimate at the Swing boundary, but they should be controlled.
A text field binding reads a model value and writes a model value. A checkbox
binding does the same. A table selection binding reads model selection and writes
model selection. The binding must prevent endless echo by comparing values and by
using origins when the distinction matters. Core values make this possible
because equal values are silent and each event carries origin. They do not make a
bad feedback design good.

\begin{lstlisting}[caption={A derived value should compute, not cause effects.}]
DerivedValue<Boolean> canDelete = DerivedValue.of(
        () -> selected.value() != null
                && !busy.value()
                && permissions.value().mayDeleteCustomers(),
        selected, busy, permissions);

selected.subscribe(event -> {
    if (event.origin() == ChangeOrigin.USER) {
        auditSelection(event.newValue());
    }
});
\end{lstlisting}

The first block computes a fact. The second block performs an effect. Keeping
them separate makes each one easier to test. The derived value test changes
selection, busy state, and permissions. The audit test changes selection with a
user origin and with a model origin. A button click is not needed to prove either
rule.

Reading several values directly is sometimes enough, but it is not always a
stable design. Suppose a Save command reads \api{dirty}, \api{problems}, and
\api{busy} at invocation time. That may be correct. Suppose a button, menu item,
keyboard shortcut, toolbar state, status text, and test all need the same
enablement. Then the enablement should be a named derived value. The distinction
is reuse and meaning. One local decision can be local code. A repeated screen
fact deserves a value.

There is no automatic transaction across several values. If code sets
\api{busy} to true and then clears \api{selectedCustomer}, observers may see both
changes separately. Often that is fine. Sometimes it is not. When a group of
facts must change as one conceptual state, introduce a single state record or a
higher-level model method that updates observers according to the model's
contract. Do not assume that three independent values will behave like one
transaction merely because they are updated in adjacent lines.

\begin{lstlisting}[caption={A combined screen mode avoids impossible states.}]
sealed interface DetailState {
    record Empty() implements DetailState {}
    record Loading(CustomerId id) implements DetailState {}
    record Ready(CustomerId id, CustomerDetails details) implements DetailState {}
    record Failed(CustomerId id, String message) implements DetailState {}
}

SimpleValue<DetailState> detailState =
        SimpleValue.of(new DetailState.Empty());
\end{lstlisting}

The combined state does not prevent every bug, but it narrows the legal space.
There is no state where the detail is both loading and failed, or ready without
an id. Derived values can expose the parts the view needs: \api{detailVisible},
\api{busyMessage}, \api{retryEnabled}, and \api{currentDetails}. The model still
has one source of truth.

Mapping and deriving are also documentation tools. The code
\api{MappedValue.of(selectedCustomer, CustomerRow::displayName)} says that the
display name follows selection. A listener named \api{selectionChanged} might do
the same thing, or it might update ten unrelated widgets. The value object has a
smaller contract. This is why derived values belong in the presenter, not hidden
inside arbitrary component factories. A generated binding can install the
relationship, but the presenter's model should tell the story.

Choose where to expose values carefully. A presenter may expose
\api{ReadableValue<Boolean>} for \api{canSave} and keep the writable dirty flag
private. A form model may expose field values through field models rather than
raw writable values. A task handle may expose busy and progress as readable
values while task control remains with a service. Exposing a writable value is a
promise that the caller may write it. If the caller should only observe, expose
the readable interface.

\begin{lstlisting}[caption={Expose readable facts, keep writable ownership local.}]
final class CustomerScreenModel {
    private final SimpleValue<Boolean> dirty = SimpleValue.of(false);
    private final DerivedValue<Boolean> canSave;

    CustomerScreenModel(ReadableValue<ProblemSet> problems,
            ReadableValue<Boolean> busy) {
        canSave = DerivedValue.of(
                () -> dirty.value()
                        && !busy.value()
                        && !problems.value().hasErrors(),
                dirty, busy, problems);
    }

    ReadableValue<Boolean> canSave() {
        return canSave;
    }

    void markDirty(ChangeOrigin origin) {
        dirty.setValue(true, origin);
    }
}
\end{lstlisting}

This pattern is especially useful with generated code. Generated bindings often
need readable facts for labels, enabled state, visible state, problems, and
tooltips. They should not receive write access unless the generated behavior is
the intended owner of a user edit. The type boundary is small, but it prevents
whole classes of accidental writes.

Origins should remain stable diagnostic ids. A custom origin created with
\api{ChangeOrigin.of} should be named as a stable fact, not as a sentence that
changes whenever a method is renamed. Good custom origins include
\api{import.customers}, \api{restore.preferences}, and \api{generated.form}.
Poor origins include \api{some button thing} or a localized message. Origins may
appear in diagnostics and tests; treat them like keys.

The event source should usually be boring. Most listeners observe one value and
already know the source. The source becomes useful when a diagnostic listener or
small adapter observes several values. For example, a presenter test may record
all changes to \api{dirty}, \api{canSave}, and \api{statusText}. The source lets
the test distinguish events without wrapping each listener manually. Production
code should not build clever dispatch systems around source identity unless the
design truly needs one.

Value graphs are not a replacement for commands. A command has identity,
resources, enablement, optional selected state, shortcuts, and invocation
behavior. A value can model the command's enabled fact or disabled reason, but
clicking Save is not setting \api{canSave} to true. The command chapter will
return to this boundary. For now, remember that values describe facts; commands
perform operations.

Value graphs are not a replacement for forms. A form field has raw text, parsed
value, validation problems, dirty state, touched state, and commit/reset policy.
A scalar value is one part of that story. A form model groups those facts and
gives them workflow meaning. If a chapter example uses a simple
\api{SimpleValue<String>} for a text field, that is a teaching simplification,
not a recommendation to replace every field model with a raw value.

Value graphs are not a replacement for observable lists. A list has insertions,
removals, replacements, permutations, filters, sorters, selected ids, and table
adapters. A scalar value may contain the selected id, the current filter string,
or the list load state. The rows themselves need collection semantics when they
change. Using a scalar list snapshot is acceptable when the entire collection is
loaded or replaced as one fact, but it is not the same as precise large-data
updates.

\begin{principlebox}{one abstraction per kind of change}
Use a value for one changing scalar fact, an observable list for collection
membership and ordering, a form model for editable field workflow, a command for
an operation, and a task handle for background work. Confusion begins when one
abstraction is asked to impersonate all the others.
\end{principlebox}

There is a practical reading order for a value graph. First read source values
and their owners. Then read derived values and ask whether their suppliers are
pure. Then read subscriptions and ask which ones perform effects. Then read
bindings and ask which side may write. Finally read close or detach methods and
make sure every owned relationship dies. This order is faster than reading from
top to bottom in a component class because it follows the architecture rather
than the sequence of constructor calls.

The same order helps during debugging. If a label is wrong, find the value bound
to the label. If that value is mapped, inspect the source. If the source is
wrong, inspect who writes it and with what origin. If the source is right but
the label is wrong, inspect the binding. If events continue after closing the
screen, inspect the subscription owner. A design that supports this debugging
path is already better than one that requires breakpoints in every listener.

\section{Lifecycle, Tests, and Migration Reviews}

Values look harmless because they are small. Their lifecycle still matters. A
\api{SimpleValue} owned by a presenter can be collected when the presenter is
collected. A \api{MappedValue} or \api{DerivedValue} owns subscriptions to its
dependencies and must be closed when the owning presenter closes. A
\api{MasterDetailValue} owns a master subscription and a cached detail. A
\api{LoadableValue} may keep a cached object and subscribers. None of these
relationships should survive the screen that made them useful.

The implementation contracts are intentionally explicit. A mapped value computes
its initial value immediately, subscribes to its source, recomputes
synchronously when the source changes, and emits only when the mapped value
changes. Closing it removes the source subscription, clears listeners, and makes
later subscriptions empty. A derived value does the same for one or more
dependencies. This is why derived and mapped values should be fields of an owner
that has a close method, not temporary objects created and forgotten in a view
factory.

\begin{lstlisting}[caption={A presenter owns the derived values it creates.}]
final class DetailPresenter implements AutoCloseable {
    private final DerivedValue<Boolean> canReload;
    private final MappedValue<DetailState, String> titleText;

    DetailPresenter(ReadableValue<DetailState> state,
            ReadableValue<Boolean> busy) {
        canReload = DerivedValue.of(
                () -> state.value() instanceof DetailState.Failed
                        && !busy.value(),
                state, busy);
        titleText = MappedValue.of(state, this::titleFor);
    }

    @Override
    public void close() {
        canReload.close();
        titleText.close();
    }
}
\end{lstlisting}

The close method is not ceremonial. It expresses ownership. Later chapters will
replace manual lists of close calls with scopes, but the responsibility remains
the same. If a derived value subscribes to a long-lived application preference
and the screen never closes it, the preference value may keep the screen graph
reachable. That is a leak with a polite API.

Loadable and detachable values add another lifecycle distinction. Closing is
about the owner being finished. Detaching is about releasing cached data while
the owner may be reused. A reusable view tab may detach detail data when hidden
and reload it when shown. A support inspector may invalidate cached diagnostics
after a model change. A wizard page may release a preview object while keeping
field state. Detach is not a value-change event by itself; refresh is the
operation that reloads and notifies if the effective value changes.

\begin{lstlisting}[caption={Detach releases cached data without pretending to edit it.}]
LoadableValue<CustomerPreview> preview =
        LoadableValue.of(() -> repository.loadPreview(currentId.value()));

CustomerPreview first = preview.value();  // load
CustomerPreview again = preview.value();  // cached

preview.detach();                         // release cache

CustomerPreview reloaded = preview.value();
\end{lstlisting}

The loader runs on the caller's thread. That sentence is easy to overlook and
important. A loadable value is not a background task. If the loader can block,
do not call it from painting or from an EDT action without a task boundary. Use
the task chapter's patterns to load in the background and publish a value on the
EDT. A loadable value is a synchronous cache abstraction; background work is a
separate concern.

Master-detail values combine lazy attachment with a master subscription. The
detail loader receives the current master, which may be null. When the detail is
attached and the master changes, the detail reloads and notifies if the
effective detail changes. When detached, a master change invalidates the stale
association but does not load until the next access. This is a precise contract
and worth explaining to readers because it avoids hidden repository calls while
a collapsed detail pane is not visible.

\begin{lstlisting}[caption={Master changes reload attached details only.}]
SimpleValue<CustomerId> selectedId = SimpleValue.empty();
MasterDetailValue<CustomerId, CustomerDetails> details =
        MasterDetailValue.of(selectedId, repository::loadDetails);

details.detach();                         // detail pane hidden
selectedId.setValue(new CustomerId(10));   // no eager load while detached

CustomerDetails visible = details.value(); // load when pane needs it
\end{lstlisting}

This pattern is not a cure for slow repositories. If \api{loadDetails} blocks,
the call to \api{details.value()} blocks. The proper architecture may be a
master value that starts a cancellable task and a detail state value that
receives results. The synchronous master-detail value is still useful for cheap
derived models, cached in-memory details, and examples where the relationship is
more important than threading.

Testing core values should be direct and small. A test for \api{SimpleValue}
sets a new value and records one event. A test for effective equality sets an
equal value and records no event. A test for origins writes with user and model
origins and asserts the difference. A test for a derived value mutates sources
and asserts the recomputed value. A test for close closes a derived value and
then mutates the source. No robot, no window, and no EDT are required.

\begin{lstlisting}[caption={Testing close semantics without Swing.}]
SimpleValue<Integer> source = SimpleValue.of(1);
DerivedValue<Integer> doubled =
        DerivedValue.of(() -> source.value() * 2, source);
List<Integer> observed = new ArrayList<>();

doubled.subscribe(event -> observed.add(event.newValue()));
doubled.close();
source.setValue(2, ChangeOrigin.MODEL);

assertThat(observed).isEmpty();
assertThat(doubled.value()).isEqualTo(2);
\end{lstlisting}

The last assertion reflects the current implementation contract: closing stops
dependency notifications; it does not erase the last computed value. A test like
this is useful because it documents what "closed" means for readers and for
future maintainers. If a later implementation changes closed-value behavior, the
test will force a deliberate decision.

Test origins as behavior, not decoration. It is not enough to assert that the
final value is correct after a table click or model load. If origins are used to
prevent feedback loops or to drive audit behavior, they are part of the
contract. A table selection binding should write user-originated changes for
user gestures. A presenter restoring selection after loading should use a model
origin. A generated binding may use a generated or binding origin according to
the binding contract. Tests should make these distinctions visible.

\begin{lstlisting}[caption={Recording origins as part of the rule.}]
SimpleValue<CustomerRow> selected = SimpleValue.empty();
List<ChangeOrigin> origins = new ArrayList<>();
selected.subscribe(event -> origins.add(event.origin()));

selected.setValue(rowFromUserClick, ChangeOrigin.USER);
selected.setValue(rowFromRestore, ChangeOrigin.MODEL);

assertThat(origins).containsExactly(ChangeOrigin.USER, ChangeOrigin.MODEL);
\end{lstlisting}

Avoid tests that assert only listener counts. A listener count may prove that an
event fired, but not that the event carried the correct old value, new value,
source, and origin. Prefer tests that record the event object. The event object
is the contract. It is also how diagnostics will explain the screen later.

Migration from plain Swing should be incremental. Do not rewrite the whole
screen into values on the first pass. Find one fact that is currently stored in
a component or private field and read in several places. Selected row is a common
candidate. Save enablement is another. Status text is another. Introduce one
source value, derive one or two presentation facts, and bind the existing
components. Run tests. Then repeat. Each step should reduce the number of places
where the fact is guessed.

\begin{lstlisting}[caption={The first migration step names selected identity.}]
// Before: listeners read table.getSelectedRow() wherever they need selection.
private final SimpleValue<CustomerId> selectedCustomerId = SimpleValue.empty();

void tableSelectionChanged() {
    CustomerId id = selectedIdFromTable();
    selectedCustomerId.setValue(id, ChangeOrigin.USER);
}

ReadableValue<Boolean> hasSelection =
        MappedValue.of(selectedCustomerId, Objects::nonNull);
\end{lstlisting}

The first version may still have a table listener. That is acceptable. The
improvement is that the listener writes a named fact instead of updating every
dependent component. Later a Corusco table binding can write the same value. The
rest of the screen does not need to know which mechanism provided selection.

Look for listener fields during migration. A class with fields named
\api{selectedRow}, \api{isDirty}, \api{currentError}, or \api{busy} that are
written by listeners and read elsewhere is asking for values. A class with
private booleans used only inside one method is not. Good migration does not
make simple local logic more abstract. It names state that already escaped.

Look for repeated enablement code. If \api{saveButton.setEnabled},
\api{saveMenuItem.setEnabled}, and a key binding all compute the same condition,
introduce \api{canSave}. If every computation is slightly different, do not
paper over the difference with a value until the policy is clear. The value
should express one rule, not average several inconsistent rules.

Look for status strings built in listeners. Status text often starts as harmless
UI feedback and becomes a cross-cutting summary of selection, load state,
validation, and task progress. A mapped or derived status value is easier to
test and localize. It is also easier to replace with a richer status object when
the screen grows.

Look for background callbacks that write components directly. A callback that
sets a label and enables a button after a repository call is mixing task result,
state mutation, and presentation. A better shape is: task result arrives at the
EDT boundary, presenter writes a state value, derived values update status and
enablement, bindings update components. This shape is slightly more vocabulary
up front and much less guesswork later.

\begin{migrationbox}{from callback mutation to model mutation}
When an asynchronous callback updates several Swing components, introduce a
single state value for the result first. Then derive component-facing facts from
that state. The callback should publish the model result, not become a second
view constructor.
\end{migrationbox}

Reviewers should be suspicious of values created inside paint, renderer, or
editor methods. A table cell renderer should not create a \api{SimpleValue} for
each cell. A custom component should not allocate derived values during painting.
Values are model objects with lifetimes; renderers and paint methods are
presentation passes. If a renderer needs derived presentation state, prepare it
in the model or descriptor layer and let the renderer read it.

Reviewers should also be suspicious of subscriptions created without visible
ownership. A call to \api{subscribe} whose result is ignored is usually a bug.
There are rare cases where a value owns the listener internally and returns a
subscription to satisfy the API, but application code should normally keep the
subscription in a scope, field, or derived object. The cleanup handle exists to
be used.

\begin{lstlisting}[caption={Ignoring a subscription hides lifecycle.}]
// Smell: who closes this listener?
statusText.subscribe(event -> statusLabel.setText(event.newValue()));

// Better: put the binding or subscription in a scope owned by the screen.
subscriptions.add(statusText.subscribe(event ->
        statusLabel.setText(event.newValue())));
\end{lstlisting}

The better listing is still not the final Corusco binding style. It is the
minimum responsible plain Java style. Later chapters introduce scopes and
binding factories to reduce this boilerplate. The ownership idea must be learned
before the helper API can be trusted.

Value design affects accessibility and help. A label's visible text, accessible
description, tooltip, help topic, and validation summary often derive from the
same field state. If those strings are all built independently inside
components, they drift. If the relevant state is modeled, each presentation
surface can derive its own text from the same fact. This is one reason values
appear early in the book: they are not merely convenience wrappers; they are the
first defense against semantic duplication.

Value design affects screenshots too. A screenshot harness can set source
values directly to produce deterministic states: no selection, selected row,
busy load, validation failure, failed detail load. If those states are hidden in
component internals, screenshot setup becomes a script of clicks and timing
assumptions. Model-owned values make visual documentation easier to generate and
easier to trust.

Finally, value design affects generated code. Annotation processors can generate
bindings, behavior plans, command adapters, and companion tests only when there
are stable names and contracts to attach to. A generated plan can bind
\api{statusText} to a label. It cannot reliably discover that three anonymous
listeners collectively mean "selected customer status." Corusco generation is
more useful when the handwritten presenter has already named the screen facts.

The review standard for this chapter is therefore practical. A reader should be
able to take a small Swing screen, identify component-owned facts, move one fact
into a value, derive two dependent facts, bind existing widgets, write a core
test, and close the subscriptions. That is enough to change the shape of the
codebase. The advanced chapters build from that move, not around it.

\section{Common Value Patterns}

The first common pattern is selected identity. Plain Swing often asks a table for
its selected view row whenever a command runs. That couples every command to the
table, the sorter, and the current selection model. A value-based screen usually
names the selection once. The table binding writes a domain id or row object.
Commands, status text, detail loading, and tests read the value. This is the
same table lesson from earlier chapters expressed as a scalar fact: the selected
thing is not the same as the selected pixel row.

\begin{lstlisting}[caption={Selection as a domain value, not a table coordinate.}]
SimpleValue<CustomerId> selectedCustomerId = SimpleValue.empty();

DerivedValue<Boolean> hasSelection = DerivedValue.of(
        () -> selectedCustomerId.value() != null,
        selectedCustomerId);

MappedValue<CustomerId, String> selectionStatus = MappedValue.of(
        selectedCustomerId,
        id -> id == null ? "No customer selected" : "Customer " + id.value());
\end{lstlisting}

This pattern also clarifies persistence. If a screen remembers selection across
refresh, it should remember a stable id, not a view row. If a generated table
binding receives a selected-id value, it can restore selection after sorting or
loading by looking for the id. If the selected object disappears, the value can
be set to null with a model origin and the rest of the screen responds. The
table is no longer the only owner of selection.

The second pattern is disabled reason. Many screens have a boolean
\api{canSave}. Users often need more than a disabled button. They need to know
why Save is disabled: no changes, validation errors, background save in
progress, missing permission, or read-only record. A boolean is enough for
enablement. A reason value is better for status text, tooltip text, help, tests,
and diagnostics.

\begin{lstlisting}[caption={Model disabled reason before adapting it to a button.}]
sealed interface SaveState {
    record Enabled() implements SaveState {}
    record Disabled(String reason) implements SaveState {}
}

DerivedValue<SaveState> saveState = DerivedValue.of(() -> {
    if (busy.value()) {
        return new SaveState.Disabled("Save is already running");
    }
    if (!dirty.value()) {
        return new SaveState.Disabled("No changes to save");
    }
    if (problems.value().hasErrors()) {
        return new SaveState.Disabled("Fix validation errors first");
    }
    return new SaveState.Enabled();
}, busy, dirty, problems);
\end{lstlisting}

The button binding may map \api{SaveState.Enabled} to \api{setEnabled(true)}.
The tooltip binding may map a disabled state to a reason. A command test may
assert that validation errors win over dirty state. A help overlay may use the
same reason. This is not overengineering when the reason is visible in several
places. It is ordinary modeling.

The third pattern is committed text versus editing text. A text field often
changes on every keystroke, but not every keystroke should rebuild a table
filter, query a repository, or mark a preference dirty. Sometimes the screen
needs two values: \api{draftFilterText} and \api{committedFilterText}. The field
binding writes the draft. Pressing Enter or clicking Apply writes the committed
value. The table filter observes the committed value. A status label may observe
both and report that a draft differs from the applied filter.

\begin{lstlisting}[caption={Draft and committed values separate typing from policy.}]
SimpleValue<String> draftFilter = SimpleValue.of("");
SimpleValue<String> committedFilter = SimpleValue.of("");

DerivedValue<Boolean> filterDirty = DerivedValue.of(
        () -> !Objects.equals(draftFilter.value(), committedFilter.value()),
        draftFilter, committedFilter);

void applyFilter() {
    committedFilter.setValue(draftFilter.value(), ChangeOrigin.USER);
}
\end{lstlisting}

This pattern prevents a common Swing problem: a document listener becomes a
search engine. The user types one character and the screen reloads. The user
types another character and the previous load races the next load. A value graph
lets the screen say when text is merely being edited and when it becomes a
screen policy.

The fourth pattern is busy as a modeled fact rather than a disabled-component
storm. Plain Swing code often disables every relevant widget at task start and
reenables it at task end. The error paths then forget one widget. A value-based
screen writes \api{busy}. Commands derive enablement from it. Overlays derive
visibility from it. Status text derives a message from it. Tests assert the
state transition once.

\begin{lstlisting}[caption={Busy state as a source for several presentation facts.}]
SimpleValue<Boolean> busy = SimpleValue.of(false);

DerivedValue<Boolean> controlsEnabled =
        DerivedValue.of(() -> !busy.value(), busy);

MappedValue<Boolean, String> busyStatus = MappedValue.of(
        busy,
        active -> active ? "Loading customers..." : "Ready");
\end{lstlisting}

A boolean busy value is sometimes too small. If the screen can load customers,
save a record, export a report, and refresh permissions independently, use a
task-state model or separate busy facts. A single global \api{busy} becomes
misleading when several operations can overlap. The pattern is not "one busy
flag everywhere"; it is "model the busy facts that the screen actually needs."

The fifth pattern is preference shadowing. Application preferences may be
long-lived, while a dialog edits a temporary copy. The dialog should not write
the global preference value on every keystroke if Cancel should discard changes.
Use dialog-local draft values and commit them explicitly. The global preference
value may be observable so open screens can react after OK. The dialog-local
values die with the dialog.

\begin{lstlisting}[caption={A preference dialog edits local values and commits.}]
SimpleValue<String> draftTheme = SimpleValue.of(preferences.theme().value());
SimpleValue<Integer> draftPageSize =
        SimpleValue.of(preferences.pageSize().value());

void commitPreferences() {
    preferences.theme().setValue(draftTheme.value(), ChangeOrigin.USER);
    preferences.pageSize().setValue(draftPageSize.value(), ChangeOrigin.USER);
}
\end{lstlisting}

This pattern is small but important for user trust. OK and Cancel semantics are
state ownership semantics. If the dialog writes global values while the user is
still browsing options, Cancel becomes a lie. Values make the two lifetimes
visible: draft lifetime and application preference lifetime.

The sixth pattern is value-backed diagnostics. A support panel can subscribe to
selected ids, load state, busy state, problem count, and command state. It can
show those facts without reaching into components. It can also record recent
origin transitions. This turns "the UI looked wrong" into a small timeline:
filter changed by user, task started by system, selection cleared by model,
problems updated by validation, save disabled because busy. That timeline is
only possible if facts are modeled before pixels.

\begin{lstlisting}[caption={A tiny diagnostic event line from a value change.}]
String describe(ValueChangeEvent<?> event) {
    return event.source().getClass().getSimpleName()
            + " "
            + event.oldValue()
            + " -> "
            + event.newValue()
            + " origin="
            + event.origin();
}
\end{lstlisting}

The source class name alone is not enough for production diagnostics, but the
shape is useful. Real diagnostics should include stable keys or field names.
Generated companions can help by assigning keys to values, commands, and
components. The core value event already contains the four facts that matter:
source, old value, new value, and origin.

The seventh pattern is the boundary between value and view lifecycle. A value
owned by a long-lived model may outlive several views. A binding from that value
to a label should not. When the view opens, it subscribes. When the view closes,
it unsubscribes. The value remains. This pattern is common for preferences,
workspace selection, logged-in user, connection state, and application theme.
It is also where memory leaks appear if subscriptions are ignored.

\begin{lstlisting}[caption={Long-lived value, short-lived view subscription.}]
final class StatusView implements AutoCloseable {
    private final Subscription subscription;

    StatusView(ReadableValue<String> applicationStatus, JLabel label) {
        label.setText(applicationStatus.value());
        subscription = applicationStatus.subscribe(event ->
                label.setText(event.newValue()));
    }

    @Override
    public void close() {
        subscription.close();
    }
}
\end{lstlisting}

This is plain Java and intentionally repetitive. Corusco binding scopes reduce
the repetition later. The important lesson is that the value's lifetime and the
view relationship's lifetime are different. A global status value is fine. A
global subscription from that value to a closed label is not.

The eighth pattern is replacing sentinel component state. Swing components often
use visible strings as state: a combo box item named "All", a label containing
"No selection", a text field containing a placeholder, a disabled button
tooltip. These strings are presentation, not durable state. A value should hold
the semantic fact, such as \api{Optional<CustomerId>} or a small filter mode.
The component adapter can render "All" or "No selection" according to resources.

\begin{lstlisting}[caption={Semantic filter mode instead of magic combo text.}]
sealed interface CustomerFilterMode {
    record All() implements CustomerFilterMode {}
    record ActiveOnly() implements CustomerFilterMode {}
    record AssignedTo(UserId userId) implements CustomerFilterMode {}
}

SimpleValue<CustomerFilterMode> filterMode =
        SimpleValue.of(new CustomerFilterMode.All());
\end{lstlisting}

This pattern matters for localization and tests. A test can assert
\api{CustomerFilterMode.All} without knowing how "All" is translated. A help
system can describe the selected mode from a resource key. A generated binding
can render the combo item without making the visible string the state.

The ninth pattern is value-backed review notes in examples. Each substantial
example should make the value graph discoverable: source values at the top,
derived values nearby, subscriptions owned by a scope, and Swing construction
after state construction. This order helps the reader. It also helps screenshots
because the harness can set source values before painting. Examples that hide
values inside anonymous listeners teach the opposite of this chapter.

The tenth pattern is knowing when to stop. A small modal confirmation dialog may
not need a value graph. A one-off label updated once after construction may not
need a \api{SimpleValue}. A local variable inside one action listener is still a
good Java local variable. Corusco is most useful where a fact crosses boundaries:
several components, several commands, several tests, several lifecycle phases, or
generated code. Use that crossing as the trigger.

These patterns repeat across the later chapters. Observable lists use selected
ids and filter text. Forms use field state, dirty state, and problem state.
Commands use enabled facts and disabled reasons. Task services publish busy and
progress facts. Bindings connect values to Swing components. Generated code
needs stable names for the same facts. The core value chapter is therefore not a
detour before the "real" library. It is the small grammar from which the later
sentences are built.

When choosing a pattern, prefer a decision made from the reader's point of view.
If the fact is written in one place and read in one place, plain Java may be
enough. If the fact is written in one place and read by several components,
commands, tests, or generated bindings, use a value. If the fact is computed
from other facts, use a mapped or derived value. If the fact is an operation,
use a command. If the fact is a collection, use an observable list. If the fact
is one part of an editable field, use a form or field model. This decision table
is simple, but it prevents the two common errors: hiding shared state in
components and using values for problems that need a richer model.

The same decision can be made by looking at failure reports. "The Save button is
enabled when it should not be" points to a derived value or command enablement.
"The row changes but the detail panel still shows the old customer" points to a
selected-id value, master-detail relationship, or stale task generation. "The
field says invalid but the OK button is still enabled" points to a problem set
and a committable derived value. "The label is translated differently from the
tooltip" points to duplicated presentation text rather than a modeled fact plus
resources. Each bug report names a missing or misplaced fact.

A value can be too small, too large, or just right. Too small: five booleans that
can represent impossible combinations. Too large: one mutable screen object
named \api{state} whose fields are changed in place and never publish effective
changes. Just right: one immutable state record for a coherent mode, or one
scalar value for a fact that truly changes independently. The art is not in
using the most values. The art is in choosing the facts a reader can verify.

Generated code benefits from this discipline because generation likes stable
edges. A generated binding can say "bind this readable value to this component
property." A generated command adapter can say "observe this readable enabled
fact." A generated test companion can say "set this source value and assert this
component state." Generation cannot rescue a screen whose state exists only as
side effects inside listeners. The handwritten value graph is therefore the
contract that generation can safely extend.

Corusco's value API is intentionally small because the difficult part is not the
method count. The difficult part is architectural honesty. Which object owns the
fact? Which code may write it? Which values derive from it? Which subscribers
perform effects? Which lifecycle closes the relationship? Which tests prove the
rule? If a chapter reader can answer those questions for a screen, the later
Corusco features will feel like helpers rather than magic.

One last review habit is useful: draw the value graph on paper for a troublesome
screen. Put writable values on the left, derived values in the middle, Swing
components and commands on the right, and subscriptions or bindings as arrows.
Mark each arrow with its owner. Mark each write with an origin. Mark each
background handoff with an EDT boundary. The diagram usually exposes the bug
before the debugger does: an unowned subscription, a component writing a fact it
should only display, a derived value doing work, or a global value pretending to
be dialog-local state.

This paper exercise also explains why the API stays small. The value graph is
not meant to impress the reader with machinery. It is meant to make ordinary
screen state inspectable. Once the graph can be read, the rest of the Corusco
stack has somewhere honest to attach.

\section{Exercises}

\begin{exercisebox}
\begin{enumerate}
\item Find a label whose text is updated from more than one listener. Name the
underlying fact and model it as a \api{ReadableValue<String>}.
\item Convert a selected-row integer into a typed selected-row value. Decide how
to represent "nothing selected".
\item Write a derived value for Save enablement using dirty state, problem
state, and busy state.
\item Add a test that records \api{ChangeOrigin.USER} and
\api{ChangeOrigin.MODEL} transitions separately.
\item Create a \api{MappedValue} for a window title. Close it and verify that
later source changes do not notify the closed mapping.
\item Identify one place where a value is not enough and a form model or command
would be the better abstraction.
\end{enumerate}
\end{exercisebox}

\textbf{Checklist.}

\begin{itemize}
\item Use one value for one changing presentation fact.
\item Keep EDT dispatch at the Swing boundary, not inside core values.
\item Treat every subscription as a resource with a visible owner.
\item Prefer mapped or derived values to listener chains that update nearby
variables.
\item Use \api{ChangeOrigin} when feedback loops, diagnostics, or tests need to
know why a value changed.
\item Model absent selection deliberately; do not let null spread without a
meaning.
\item Close derived, mapped, master-detail, and loadable values when their
lifetime ends.
\item Test value rules without constructing Swing components.
\end{itemize}
